\section{The LLM Interpretation Layer}\label{sec:llm}

Statistical outputs --- cluster centroids, support counts, metric tables --- are
not directly actionable for a clinician or a manager. HospitAI attaches a large
language model (LLM), accessed through a hosted API, to every analytical stage,
in three distinct roles. In all three, the LLM is constrained to operate over
\emph{computed statistics or validated queries}, never over free-form clinical
advice, and its textual output is explicitly a description of the data rather
than a medical recommendation.

\subsection{Cluster summarization}
For each cluster the system assembles an aggregated profile --- size,
key averages (age, severity, counts), disease prevalences, dominant care-plan
mix, and the top conditions, medications and procedures --- and asks the LLM to
produce a short structured summary: a plain-language title, a one- or two-sentence
description, the most distinguishing insight (with concrete numbers), a risk
level and a recommended intervention. The structured response is parsed and
cached on the cluster record. This converts the numerical clustering output into
something a non-specialist can read at a glance.

\subsection{Sequence interpretation}
On demand, a mined trajectory can be explained. The system gathers the
demographics and comorbidity rates of the patients that follow the pattern and
asks the LLM to produce a clinical interpretation of the trajectory (the likely
diagnosis--workup--management story), a layperson summary, risk signals and
recommendations. As with clusters, the result is cached on the pattern record so
that the (cost-bearing) model call happens at most once per pattern, with
row-level locking to serialize concurrent requests.

\subsection{Guarded text-to-SQL}
The third role lets users query the database in natural language. The agent is
implemented as a state machine (using LangGraph) with a
classify $\rightarrow$ generate $\rightarrow$ validate $\rightarrow$ execute
flow. An input is first classified as a database question, a general medical
question or off-topic. Database questions are translated into a read-only
\texttt{SELECT} grounded in the real schema; a separate validation step then
audits the query --- rejecting any non-\texttt{SELECT} statement, verifying that
referenced tables and columns exist, checking that joins use documented foreign
keys, and enforcing role-based scope (a doctor's queries are restricted to their
own patients; aggregate roles are unrestricted; the patient role is blocked).
Validation or execution failures feed the error back to the generation step for a
bounded number of retries. The natural-language answer is returned together with
the generated SQL, so a user can audit exactly what ran against the data without
needing to understand it. This design keeps the convenience of natural-language
access while preserving safety and grounding.
